% GENERATED FILE. Edit canonical lessons, then run npm run generate:books.
% canonical-source-hash: fnv1a64:5d9624deebdd8c22
% canonical-lessons: UR-C175-ras-malai, UR-C175-faluda, UR-C175-qulfi, UR-C175-bater, UR-C175-fakhta

\chapter{Rasmalai, Falooda, Kulfi, Quail, Dove}
\label{ch:ur-rasmalai-falooda-kulfi-quail-dove}
\hlchaptermodality{175}

\begin{chapteropening}
\textbf{By the end of this chapter:} \emph{I can say rasmalai, falooda, kulfi, a quail and a dove.}

Complete the last lesson of chapter 175: I can say rasmalai, falooda, kulfi, a quail and a dove.
\end{chapteropening}

\section[\texorpdfstring{\ur{رس} \ur{ملائی} (ras malāī)}{رس ملائی (ras malāī)}]{\ur{رس} \ur{ملائی} (ras malāī) — rasmalai}

\label{lesson:UR-C175-ras-malai}



\begin{quote}
Before the new one: say the Urdu for kidney beans, then the Urdu for a fig.
\end{quote}

\subsection*{You'll want to know: \ur{رس} \ur{ملائی}}
\textbf{\ur{رس} \ur{ملائی}} — \emph{ras malāī} — \textquotedblleft{}rasmalai\textquotedblright{}.

\begin{cousinweb}
\textbf{\ur{رس} \ur{ملائی}} (\emph{ras malāī}) is \textquotedblleft{}cream in syrup\textquotedblright{}. \emph{ras}, \textquotedblleft{}juice, syrup\textquotedblright{}, comes from Sanskrit \emph{rasa}, \textquotedblleft{}juice\textquotedblright{} (CDIAL 10650). The origin of \emph{malāī}, \textquotedblleft{}cream\textquotedblright{}, is not settled; CDIAL has no entry for it.
\end{cousinweb}

\begin{grammarlens}[title={What you've built}]
One more word for this chapter.
\end{grammarlens}

\subsection*{Your turn}
\begin{itemize}
\raggedright
  \item \emph{Say it:} \emph{ras malāī}
  \item \emph{Say it:} \emph{ras malāī}, once more
  \item \emph{Say it:} say \emph{ras malāī}
  \item \emph{Recall:} say the Urdu for kidney beans, then the Urdu for a fig, then say \emph{ras malāī} again
\end{itemize}

\begin{tcolorbox}[breakable,colback=teal!4,colframe=teal!35!black,title={Before you move on}]
What does \ur{رس} \ur{ملائی} mean? (\textquotedblleft{}Rasmalai\textquotedblright{}.) Say it once more.
\end{tcolorbox}
\section[\texorpdfstring{\ur{فالودہ} (fālūdā)}{فالودہ (fālūdā)}]{\ur{فالودہ} (fālūdā) — falooda}

\label{lesson:UR-C175-faluda}



\begin{quote}
Before the new one: say the Urdu for a fig, then the Urdu for rasmalai.
\end{quote}

\subsection*{You'll want to know: \ur{فالودہ}}
\textbf{\ur{فالودہ}} — \emph{fālūdā} — \textquotedblleft{}falooda\textquotedblright{}.

\begin{cousinweb}
\textbf{\ur{فالودہ}} (\emph{fālūdā}) is usually traced to Persian \emph{pālūda}, also \emph{fālūda}, \textquotedblleft{}strained, refined\textquotedblright{}, from \emph{pālūdan}, \textquotedblleft{}to strain\textquotedblright{}.
\end{cousinweb}

\begin{grammarlens}[title={What you've built}]
One more word for this chapter.
\end{grammarlens}

\subsection*{Your turn}
\begin{itemize}
\raggedright
  \item \emph{Say it:} \emph{fālūdā}
  \item \emph{Say it:} \emph{fālūdā}, once more
  \item \emph{Say it:} say \emph{fālūdā}
  \item \emph{Recall:} say the Urdu for a fig, then the Urdu for rasmalai, then say \emph{fālūdā} again
\end{itemize}

\begin{tcolorbox}[breakable,colback=teal!4,colframe=teal!35!black,title={Before you move on}]
What does \ur{فالودہ} mean? (\textquotedblleft{}Falooda\textquotedblright{}.) Say it once more.
\end{tcolorbox}
\section[\texorpdfstring{\ur{قلفی} (qulfī)}{قلفی (qulfī)}]{\ur{قلفی} (qulfī) — kulfi}

\label{lesson:UR-C175-qulfi}



\begin{quote}
Before the new one: say the Urdu for rasmalai, then the Urdu for falooda.
\end{quote}

\subsection*{You'll want to know: \ur{قلفی}}
\textbf{\ur{قلفی}} — \emph{qulfī} — \textquotedblleft{}kulfi\textquotedblright{}.

\begin{cousinweb}
\textbf{\ur{قلفی}} (\emph{qulfī}) is usually linked with \emph{qulf}, \textquotedblleft{}a lock\textquotedblright{}, the Urdu and Persian form of Arabic \emph{qufl} (Lane, root \emph{q-f-l}): the sweet is frozen in a sealed mould.
\end{cousinweb}

\begin{grammarlens}[title={What you've built}]
One more word for this chapter.
\end{grammarlens}

\subsection*{Your turn}
\begin{itemize}
\raggedright
  \item \emph{Say it:} \emph{qulfī}
  \item \emph{Say it:} \emph{qulfī}, once more
  \item \emph{Say it:} say \emph{qulfī}
  \item \emph{Recall:} say the Urdu for rasmalai, then the Urdu for falooda, then say \emph{qulfī} again
\end{itemize}

\begin{tcolorbox}[breakable,colback=teal!4,colframe=teal!35!black,title={Before you move on}]
What does \ur{قلفی} mean? (\textquotedblleft{}Kulfi\textquotedblright{}.) Say it once more.
\end{tcolorbox}
\section[\texorpdfstring{\ur{بٹیر} (baṭer)}{بٹیر (baṭer)}]{\ur{بٹیر} (baṭer) — a quail}

\label{lesson:UR-C175-bater}



\begin{quote}
Before the new one: say the Urdu for falooda, then the Urdu for kulfi.
\end{quote}

\subsection*{You'll want to know: \ur{بٹیر}}
\textbf{\ur{بٹیر}} — \emph{baṭer} — \textquotedblleft{}a quail\textquotedblright{}.

\begin{cousinweb}
\textbf{\ur{بٹیر}} (\emph{baṭer}) is a tadbhava, a word inherited and reshaped by sound change. It comes from \emph{*vartakara}, \textquotedblleft{}making turns (of the quail)\textquotedblright{}, a form CDIAL reconstructs from its descendants (CDIAL 11350). CDIAL suggests, with a query, that it is a reshaping of Sanskrit \emph{vartikā}, \textquotedblleft{}quail\textquotedblright{}.
\end{cousinweb}

\begin{grammarlens}[title={What you've built}]
One more word for this chapter.
\end{grammarlens}

\subsection*{Your turn}
\begin{itemize}
\raggedright
  \item \emph{Say it:} \emph{baṭer}
  \item \emph{Say it:} \emph{baṭer}, once more
  \item \emph{Say it:} say \emph{baṭer}
  \item \emph{Recall:} say the Urdu for falooda, then the Urdu for kulfi, then say \emph{baṭer} again
\end{itemize}

\begin{tcolorbox}[breakable,colback=teal!4,colframe=teal!35!black,title={Before you move on}]
What does \ur{بٹیر} mean? (\textquotedblleft{}A quail\textquotedblright{}.) Say it once more.
\end{tcolorbox}
\section[\texorpdfstring{\ur{فاختہ} (fākhtā)}{فاختہ (fākhtā)}]{\ur{فاختہ} (fākhtā) — a dove}

\label{lesson:UR-C175-fakhta}



\begin{quote}
Before the new one: say the Urdu for kulfi, then the Urdu for a quail.
\end{quote}

\subsection*{You'll want to know: \ur{فاختہ}}
\textbf{\ur{فاختہ}} — \emph{fākhtā} — \textquotedblleft{}a dove\textquotedblright{}.

\begin{cousinweb}
\textbf{\ur{فاختہ}} (\emph{fākhtā}) is a loan, usually traced to Persian \emph{fākhta}, \textquotedblleft{}ring-dove\textquotedblright{}.
\end{cousinweb}

\begin{grammarlens}[title={What you've built}]
One more word for this chapter.
\end{grammarlens}

\subsection*{Your turn}
\begin{itemize}
\raggedright
  \item \emph{Say it:} \emph{fākhtā}
  \item \emph{Say it:} \emph{fākhtā}, once more
  \item \emph{Say it:} say \emph{fākhtā}
  \item \emph{Recall:} say the Urdu for kulfi, then the Urdu for a quail, then say \emph{fākhtā} again
\end{itemize}

\begin{tcolorbox}[breakable,colback=teal!4,colframe=teal!35!black,title={Before you move on}]
What does \ur{فاختہ} mean? (\textquotedblleft{}A dove\textquotedblright{}.) Say it once more.
\end{tcolorbox}
